% Owner: Huy (Phase 1--2). Lab items 3.2.1, 3.2.3, 3.3.2; docs/application/PROJECT.md §4.
% Paper facts: PROJECT.md §4.2 (checked against the paper). Code facts: official repo at
% commit 235f89e, read from experiments/reproduction/SimCSE (simcse/models.py, train.py).
\section{Proposed Method}
\label{sec:method}

This section first analyses SimCSE~\cite{gao-etal-2021-simcse}: the problem it
solves, its two training settings, and why the objective yields good sentence
embeddings (Sections~\ref{sec:method-problem}--\ref{sec:method-theory}). It then
follows the official code (Section~\ref{sec:method-code}) and our own
implementation of the loss (Section~\ref{sec:method-infonce}). Finally, it
describes how we use a SimCSE encoder in a retrieval pipeline for Vietnamese
academic regulations (Section~\ref{sec:method-pipeline}).

\subsection{Problem and Key Idea}
\label{sec:method-problem}

A sentence encoder $f$ maps a sentence $x$ to a vector
$\mathbf{h}=f(x)\in\mathbb{R}^d$. The goal is that the cosine similarity
\begin{equation}
\mathrm{sim}(\mathbf{h}_1,\mathbf{h}_2)
  = \frac{\mathbf{h}_1^\top\mathbf{h}_2}{\lVert\mathbf{h}_1\rVert\,\lVert\mathbf{h}_2\rVert}
\end{equation}
is high for sentences with similar meaning and low otherwise. Without further
training, embeddings taken from pretrained language models such as
BERT~\cite{devlin-etal-2019-bert} perform poorly on this goal. On the seven STS
tasks, averaging the first and last layers of BERT-base gives an average
Spearman correlation of only 56.70~\cite{gao-etal-2021-simcse}.

SimCSE fine-tunes all weights of a pretrained BERT or
RoBERTa~\cite{liu2019robertarobustlyoptimizedbert} encoder with a contrastive
objective. Each sentence $x_i$ in a batch is paired with a positive $x_i^+$,
and the other sentences in the same batch act as negatives. The two settings
differ only in how positives (and, optionally, hard negatives) are obtained.
\emph{Unsupervised} SimCSE uses the same sentence twice and relies on dropout
noise. \emph{Supervised} SimCSE uses sentence pairs from natural language
inference (NLI) data.

In both settings the loss updates the encoder itself. Its numerator rewards
similarity to the positive, while the denominator makes that positive
compete with other candidates. Section~\ref{sec:related-embeddings}
compares the supervision and geometry of earlier sentence encoders.

\subsection{Contrastive Objective}
\label{sec:method-objective}

Given a batch of $N$ pairs $\{(x_i,x_i^+)\}_{i=1}^N$ with embeddings
$\mathbf{h}_i$ and $\mathbf{h}_i^+$, the loss for pair $i$ is the InfoNCE loss
with in-batch negatives~\cite{pmlr-v119-chen20j} (Eq.~1
of~\cite{gao-etal-2021-simcse}):
\begin{equation}
\ell_i = -\log
  \frac{e^{\mathrm{sim}(\mathbf{h}_i,\mathbf{h}_i^+)/\tau}}
       {\sum_{j=1}^{N} e^{\mathrm{sim}(\mathbf{h}_i,\mathbf{h}_j^+)/\tau}},
\label{eq:infonce}
\end{equation}
where $\tau$ is a temperature ($\tau=0.05$ by default in the official code).
Equation~\eqref{eq:infonce} is an $N$-way classification problem. Row $i$ of
the $N\times N$ matrix of scaled similarities holds the logits, and the correct
class is the diagonal entry $i$. The batch loss is the mean of $\ell_i$, which
is a standard cross-entropy with labels $0,1,\dots,N-1$. A small $\tau$
sharpens the softmax, so that negatives that are already close to $\mathbf{h}_i$
dominate the gradient.

\subsection{Unsupervised SimCSE: Dropout as Data Augmentation}
\label{sec:method-unsup}

In the unsupervised setting the positive is the sentence itself,
$x_i^+=x_i$. The sentence is encoded twice with independent dropout masks
$z_i$ and $z_i'$ on the standard Transformer dropout ($p=0.1$). This gives two
slightly different embeddings $\mathbf{h}_i^{z_i}$ and $\mathbf{h}_i^{z_i'}$,
and the loss becomes (Eq.~4 of~\cite{gao-etal-2021-simcse})
\begin{equation}
\ell_i = -\log
  \frac{e^{\mathrm{sim}(\mathbf{h}_i^{z_i},\mathbf{h}_i^{z_i'})/\tau}}
       {\sum_{j=1}^{N} e^{\mathrm{sim}(\mathbf{h}_i^{z_i},\mathbf{h}_j^{z_j'})/\tau}}.
\label{eq:unsup}
\end{equation}
Dropout therefore acts as a minimal data augmentation. It changes the hidden
representation but never the meaning of the input. Discrete augmentations, in
contrast, can change the meaning: cropping, word deletion, synonym replacement
and masked-token replacement all score lower than dropout on the STS-B
development set (Table~1 of~\cite{gao-etal-2021-simcse}). The noise level
matters. With $p=0.1$ the STS-B development score is 82.5. Removing dropout
($p=0$) lowers it to 71.1, and using the same mask for both copies
(``fixed 0.1'') lowers it to 43.6 (Table~3 of~\cite{gao-etal-2021-simcse}).
In both failure cases the two views are identical, so the positive pair carries
no learning signal. The model can then lower the loss only by pushing the
negatives away, which degrades the representation.

\subsection{Supervised SimCSE: NLI Pairs and Hard Negatives}
\label{sec:method-sup}

The supervised setting uses SNLI~\cite{bowman-etal-2015-large} and
MNLI~\cite{williams-etal-2018-broad}. For a premise $x_i$, its entailment
hypothesis is the positive $x_i^+$ and its contradiction hypothesis is a
\emph{hard negative} $x_i^-$: a sentence that is topically close but
contradicts the premise. All $2N$ hypotheses in the batch serve as candidates.
The training objective is Eq.~5 of~\cite{gao-etal-2021-simcse}. Its ablation
in Section~6.3 of~\cite{gao-etal-2021-simcse} adds a weight $\alpha$ on each
anchor's own hard negative (Eq.~8 there), and we write the loss in this
weighted form:
\begin{equation}
\ell_i = -\log
  \frac{e^{\mathrm{sim}(\mathbf{h}_i,\mathbf{h}_i^+)/\tau}}
       {\sum_{j=1}^{N}\left( e^{\mathrm{sim}(\mathbf{h}_i,\mathbf{h}_j^+)/\tau}
        + \alpha^{\delta_{ij}}\, e^{\mathrm{sim}(\mathbf{h}_i,\mathbf{h}_j^-)/\tau}\right)}.
\label{eq:sup}
\end{equation}
Here $\delta_{ij}$ is 1 if $i=j$ and 0 otherwise, so the weight $\alpha$
applies only to the sentence's own hard negative, and
$\alpha=1$ gives Eq.~5 itself. Adding hard negatives raises the STS-B
development score from 84.9 to 86.2 (Table~4 of~\cite{gao-etal-2021-simcse}).
% Published ablation: ../2021.emnlp-main.552.pdf (repo root), Table 7.
With contradiction hard negatives, $\alpha=1$ and $\alpha=2$ tie at the
reported precision in Table~7 of~\cite{gao-etal-2021-simcse}. In the code, the weight is passed as
$w=\log\alpha$ (argument \texttt{hard\_negative\_weight}, default $0$) and
added to the corresponding logit, because
$e^{s/\tau+w}=\alpha\,e^{s/\tau}$. Adding to a logit fits directly into the
cross-entropy implementation. An optional auxiliary masked-language-modelling
loss also exists in the code but is off by default.

\subsection{Worked Example}
\label{sec:method-example}

% The cosines are illustrative values chosen by us, not model outputs.
% tests/test_worked_example.py recomputes every loss and probability below and
% checks that each printed value appears in this file.
We illustrate both settings with the sentences of Fig.~1
of~\cite{gao-etal-2021-simcse}. The cosine values are illustrative values
chosen by us, not outputs of a trained model. The losses come from our own
implementation (Section~\ref{sec:method-infonce}), and
\path{tests/test_worked_example.py} recomputes every number below.

\emph{Unsupervised.} A batch holds three sentences: $x_1$~=~``Two dogs are
running'', $x_2$~=~``A man surfing on the sea'' and $x_3$~=~``A kid is on a
skateboard''. Each sentence is encoded twice with independent dropout masks
(in the code, both copies sit in one batch; Section~\ref{sec:method-code}).
Dropout switches off different hidden units in the two copies, so $x_1$ gets
two slightly different embeddings, $\mathbf{h}_1^{z_1}$ and
$\mathbf{h}_1^{z_1'}$. They form the positive pair. The second copies of the
other sentences, $\mathbf{h}_2^{z_2'}$ and $\mathbf{h}_3^{z_3'}$, are the
negatives of $x_1$. Suppose the cosines from $\mathbf{h}_1^{z_1}$ to
$\mathbf{h}_1^{z_1'}$, $\mathbf{h}_2^{z_2'}$ and $\mathbf{h}_3^{z_3'}$ are
$0.80$, $0.70$ and $0.60$. We set the negatives high on purpose, as for a
pretrained encoder whose embeddings are anisotropic, so that even unrelated
sentences are close (Section~\ref{sec:method-theory}). With $\tau=0.05$ the
logits are $16$, $14$ and $12$. The softmax gives the positive a probability
of $0.87$, and Eq.~\eqref{eq:unsup} gives $\ell_1=0.143$. With $\tau=1$ the
same cosines give the positive a probability of only $0.37$, close to the
uniform $1/3$, and $\ell_1=1.002$. A small $\tau$ therefore sharpens the
softmax, so that small cosine gaps decide which candidate wins. Minimising
$\ell_1$ pulls the two copies of $x_1$ together and pushes the other sentences
away.

\emph{Supervised.} Each premise now comes with its entailment and
contradiction hypotheses from the same figure. For $x_1$, the positive is
``There are animals outdoors'' and the hard negative is ``The pets are sitting
on a couch''. The other hypotheses in the batch are further negatives: ``There
is a man'', ``A kid is skateboarding'', ``The man wears a business suit'' and
``A kid is inside the house''.\footnote{Fig.~1 of~\cite{gao-etal-2021-simcse}
prints ``kit'' instead of ``kid'' in this sentence; we correct the typo.}
Row~1 of the logit matrix therefore has $2N=6$ entries. We now assume a later
stage of training, when unrelated sentences are already well separated. Let
the cosines be $0.80$ for the positive, $0.40$ and $0.45$ for the other
entailments, $0.75$ for the anchor's own hard negative, and $0.30$ and $0.35$
for the other contradictions. With $\tau=0.05$ and $\alpha=1$,
Eq.~\eqref{eq:sup} gives $\ell_1=0.314$. The positive gets probability $0.73$
and the hard negative $0.27$, while the four other negatives together get less
than $0.01$. Without the contradictions, the same row would give
$\ell_1=0.0012$. Almost all of the loss for $x_1$ thus comes from its hard
negative, which is topically close but contradicts the premise. With $\alpha=2$
($w=\log 2$) the loss rises to $0.552$. The gradient of $\ell_1$ with respect
to a negative's logit equals that negative's softmax probability, and for the
hard negative this probability rises from $0.27$ to $0.42$. A larger $\alpha$
therefore pushes the hard negative away more strongly.

\subsection{Pooling, Training and Inference}
\label{sec:method-pooling}

The sentence embedding is the final hidden state of the \texttt{[CLS]} token.
During training, a one-layer MLP (a linear layer followed by $\tanh$) is put on
top of it. For unsupervised SimCSE the MLP is used only during training and
removed at test time. Supervised SimCSE keeps the MLP at test time
(Table~6 of~\cite{gao-etal-2021-simcse}). Unsupervised models are trained on
$10^6$ sentences sampled from English Wikipedia for one epoch (BERT-base: batch
size 64, learning rate $3\times10^{-5}$). Supervised models are trained on the
combined SNLI and MNLI data for three epochs (batch size 512, learning rate
$5\times10^{-5}$). In both settings the STS-B development set is evaluated
periodically and the best checkpoint is kept (Appendix~A
of~\cite{gao-etal-2021-simcse}).

At inference time dropout is off. A sentence is encoded once, pooled as above
and compared by cosine similarity. STS evaluation uses this cosine directly,
with no trained regressor, and reports Spearman correlation
(Section~\ref{sec:setup}).

\subsection{Why the Objective Works: Alignment, Uniformity and Anisotropy}
\label{sec:method-theory}

Wang and Isola~\cite{pmlr-v119-wang20k} describe contrastive representations by
two properties of the normalised embedding $f(x)$. \emph{Alignment} measures how
close positive pairs are,
\begin{equation}
\ell_{\mathrm{align}} \triangleq
  \mathop{\mathbb{E}}_{(x,x^+)\sim p_{\mathrm{pos}}}
  \lVert f(x)-f(x^+)\rVert^2 ,
\end{equation}
and \emph{uniformity} measures how evenly embeddings spread over the unit
hypersphere,
\begin{equation}
\ell_{\mathrm{uniform}} \triangleq \log
  \mathop{\mathbb{E}}_{x,y \overset{\mathrm{i.i.d.}}{\sim} p_{\mathrm{data}}}
  e^{-2\lVert f(x)-f(y)\rVert^2}.
\end{equation}
Lower is better for both. The numerator of Eq.~\eqref{eq:infonce} rewards
alignment and the denominator rewards uniformity. Gao et al.\ track both
quantities during training (Fig.~2 of~\cite{gao-etal-2021-simcse}).
Unsupervised SimCSE steadily improves uniformity while keeping alignment
stable, whereas without dropout noise or with a fixed mask, alignment
deteriorates.

Uniformity is closely linked to \emph{anisotropy}. Embeddings from pretrained
language models occupy a narrow cone of the vector
space~\cite{ethayarajh-2019-contextual}, so even unrelated sentences have high
cosine similarity. For a uniform sample of $m$ normalised embeddings,
let $G=\mathbf{W}\mathbf{W}^\top$ be their Gram matrix. Jensen's inequality
shows that the asymptotic negative term of the contrastive loss is at least
the mean of the entries of $G$ divided by $\tau$
(Eq.~7 of~\cite{gao-etal-2021-simcse}). When all entries of $G$ are
positive, their sum also bounds its largest eigenvalue from above, while
its trace stays fixed at $m$. This motivates the paper's spectral argument.
Appendix~F of~\cite{gao-etal-2021-simcse} shows a flatter singular value
spectrum for SimCSE on the sampled STS-B sentences.

\subsection{Execution Flow of the Official Code}
\label{sec:method-code}

We read the official implementation at commit
\texttt{235f89e}.\footnote{\url{https://github.com/princeton-nlp/SimCSE/tree/235f89e0244b23673b69fbbc4a1cd9868c5a5c24}}
One training step proceeds as follows.
\begin{enumerate}
\item \emph{Data loading} (\texttt{train.py}). A \texttt{.txt} file has one
column, which is used as both the sentence and its positive. A \texttt{.csv}
file with two columns gives (sentence, positive) pairs, and one with three
columns adds a hard negative. \texttt{prepare\_features} tokenises all columns
(up to 32 tokens in the example scripts) and groups the \texttt{num\_sent}
encodings of each example together. The data collator pads the batch and
reshapes it to a tensor of shape $(N, \texttt{num\_sent}, L)$.
\item \emph{Encoding} (\texttt{cl\_forward} in \texttt{simcse/models.py}). The
tensor is flattened to $(N\cdot\texttt{num\_sent}, L)$ and encoded in a single
batched forward pass. The two copies of a sentence are separate rows of this
batch, so in training mode they receive independent dropout masks. This is how
Eq.~\eqref{eq:unsup} is implemented without running the encoder twice.
\item \emph{Pooling.} The \texttt{Pooler} takes the \texttt{[CLS]} vector. With
\texttt{pooler\_type=cls}, the MLP is applied on top. The output is reshaped to
$(N,\texttt{num\_sent},d)$ and split into $\mathbf{z}_1$, $\mathbf{z}_2$ and,
if present, $\mathbf{z}_3$.
\item \emph{Multi-GPU batches.} Under distributed training,
\texttt{all\_gather} collects the embeddings from all GPUs. The local slice is
put back so that gradients still flow through it. The effective batch is
therefore the sum over GPUs, which is how the example script reaches the
supervised batch size of 512 with four GPUs of 128 each.
\item \emph{Loss.} The \texttt{Similarity} module computes cosine divided by
$\tau$ between every pair in $\mathbf{z}_1\times\mathbf{z}_2$, giving an
$N\times N$ logit matrix. With hard negatives, the $\mathbf{z}_1\times\mathbf{z}_3$
block is concatenated, giving $N\times 2N$, and $w$ is added at column $N+i$ of
row $i$. A cross-entropy loss with labels $0,\dots,N-1$ then implements
Eqs.~\eqref{eq:infonce}--\eqref{eq:sup}.
\item \emph{Model selection} (\texttt{simcse/trainers.py}). During training,
\texttt{CLTrainer.evaluate} runs SentEval on the STS-B and SICK-R development
sets and keeps the checkpoint with the best STS-B Spearman correlation.
\item \emph{Output.} \texttt{simcse\_to\_huggingface.py} renames the
\texttt{mlp} weights to \texttt{pooler} and removes the encoder prefix. The MLP
has the same structure as the BERT pooler (linear layer and $\tanh$), so in the
converted checkpoint the standard \texttt{pooler\_output} equals the
training-time MLP output.
\end{enumerate}
At inference, \texttt{sentemb\_forward} returns the pooled vector, and
\texttt{evaluation.py} scores STS with the pooling chosen by
\texttt{-{}-pooler}. \texttt{cls} gives \texttt{pooler\_output}, used for
supervised models, while \texttt{cls\_before\_pooler} gives the raw
\texttt{[CLS]} vector, used for unsupervised models.

We note two differences between this commit and the paper. First, the commit
has no \texttt{-{}-mlp\_only\_train} option. ``MLP only during training'' is
obtained by training with \texttt{cls} and evaluating with
\texttt{cls\_before\_pooler}. Second, the example scripts evaluate every 125
steps, whereas Appendix~A of the paper states 250.

\subsection{Independent Implementation of the Loss}
\label{sec:method-infonce}

To check our understanding of Eqs.~\eqref{eq:infonce}--\eqref{eq:sup}, we
implement the loss ourselves in \path{src/simcse_lab/model/infonce.py}. It
covers the unsupervised case and the supervised case with a log-weighted hard
negative. Unlike \texttt{cl\_forward}, it normalises the embeddings, takes
their dot products, subtracts the row maximum before dividing by $\tau$,
and adds $w$ to each anchor's own hard-negative logit. The negative diagonal
of the row-wise log-softmax gives the loss for each positive pair. Centering
first preserves the contribution of tied candidates at low temperatures.
It takes the already gathered batch, so multi-GPU gathering is out of scope.

The implementation is tested on random embeddings against two references: a
naive loop that evaluates Eqs.~\eqref{eq:unsup} and~\eqref{eq:sup} term by term
with $\alpha=e^{w}$, and the loss computation copied verbatim from
\texttt{cl\_forward} (\path{tests/infonce_reference.py}). The acceptance
criterion is an absolute difference below $10^{-6}$ in float64.
% Numbers below: experiments/reproduction/results/R4-2026-11-24/metrics.json (seed 0,
% torch 2.9.1 on CPU).
The test grid (seed 0) covers batch sizes $N\in\{1,2,8,32\}$, plus
$N\in\{128,512\}$ against \texttt{cl\_forward} only, embedding sizes 16 and
768, $\tau\in\{0.01,0.05,0.1,1\}$ and
$w\in\{0,1,-0.5,\log 0.5,\log 2\}$. Positives and hard negatives are noisy
copies of the anchor, so that the hard negatives compete with the positive.
All 689 test cases pass. In float64, the largest difference is
$8.7\times10^{-15}$ from the naive loop. Against \texttt{cl\_forward} it is
$3.6\times10^{-15}$ when $w$ is exactly representable in float32, and
$1.4\times10^{-9}$ for $w=\log 0.5$ and $w=\log 2$. The larger gap comes from
\texttt{cl\_forward} building its weight matrix in float32: with that matrix
built in float64 instead, the gap falls to $5.3\times10^{-15}$. On one batch
($N=8$, $d=32$, $\tau=0.05$, $w\in\{0,\log 2\}$), the gradients with respect
to all three inputs also agree within $10^{-6}$. In float32 the two
implementations round differently, both in the cosine logits and in the
cross-entropy step, and the loss gap reaches $9.5\times10^{-7}$. The float32
check therefore uses a tolerance of $5\times10^{-5}$.

\subsection{Retrieval Pipeline for Vietnamese Regulations}
\label{sec:method-pipeline}

The application does not train a new encoder. It applies a community checkpoint
trained following the SimCSE recipe on PhoBERT-base~\cite{nguyen-tuan-nguyen-2020-phobert}
(\texttt{VoVanPhuc/sup-SimCSE-VietNamese-phobert-base}).\footnote{\url{https://hf.co/VoVanPhuc/sup-SimCSE-VietNamese-phobert-base}}
Its model card does not state the training data, so its scores are not results
of the SimCSE paper. The following design specifies offline indexing and online
query processing for the planned application.

% Corpus details: data/d2/README.md and data/d2/documents.yaml (decision D-17).
\textbf{Indexing.} The assistant indexes the undergraduate training
regulations of our university: the 2021 regulation (Decision 1175/QĐ-KHTN, 24
articles) and the 2026 regulation that replaces it (Decision 3087/QĐ-KHTN, 27
articles). Both versions are kept because the 2026 text applies to cohorts
admitted from the 2026--2027 academic year, while for earlier cohorts it lets
the university apply the new rules that are more favourable to students. For the evaluation on
ViRHE4QA, a separate index holds the distinct passages of that dataset, most of
which are single articles (\emph{Điều}). Each article is split into its
clauses (\emph{Khoản}).
% Window limit and overlap: src/simcse_lab/retrieval/ingest.py and src/simcse_lab/model/encoders.py.
Each chunk fits the 256-token input budget,
including special tokens. Longer clauses are split at whitespace with up to
24 overlapping whitespace-delimited units between consecutive windows.
Every chunk keeps its document, article number and title, clause and character
offsets, so the interface can quote the indexed text exactly. For our
university's regulations, each document also records its decision number,
issue date and scope of application, which the interface displays. For
ViRHE4QA, we recover the document title by matching passage text to the
authors' workbook. The title alone does not establish the applicable edition.
The 2026 HCMUS regulation was available to us only as a
scanned PDF. Its indexed text is a transcription of the scan made with ChatGPT,
which we first cleaned in two passes and compared with the scan's OCR layer.
We then checked discrepancies against images of the signed regulation and
corrected substantive differences, including the threshold for dismissal
after consecutive academic warnings. The interface identifies the indexed
text as a transcription. Text is normalised to
Unicode NFC, tone-mark placement is unified (e.g.\ \emph{khoá} and \emph{khóa}),
and words are segmented with pyvi, which the PhoBERT vocabulary requires. The
same function is applied to queries, so that index and queries are always
segmented in the same way. We then build a BM25 index~\cite{Robertson_2009}
over lower-cased tokens and a dense index of L2-normalised SimCSE embeddings
(\texttt{pooler\_output}). We compute each exact cosine score as the inner
product of the L2-normalised query and chunk vectors in NumPy.

\textbf{Querying.} For a question $q$, each chunk $c$ receives a BM25 score and
a cosine score. Because the gold labels and the displayed results are at the
article level, chunk scores are aggregated to an article $a$ by taking the
maximum over its chunks, following the MaxP aggregation rule~\cite{Dai_2019}:
$s(q,a)=\max_{c\in a} s(q,c)$. For each query, the BM25 and dense article
scores are min--max normalised to $[0,1]$ separately, giving $\hat{s}_{\mathrm{bm25}}$
and $\hat{s}_{\mathrm{dense}}$, and are fused linearly:
\begin{equation}
s_{\mathrm{hyb}}(q,a) = \lambda\,\hat{s}_{\mathrm{bm25}}(q,a)
  + (1-\lambda)\,\hat{s}_{\mathrm{dense}}(q,a).
\label{eq:fusion}
\end{equation}
BM25 captures exact terms, numbers and abbreviations, and the encoder captures
paraphrases. The weight $\lambda\in[0,1]$ is the only fitted parameter of the
ranking. We selected $\lambda=0.6$ on D1-dev (Section~\ref{sec:setup}).
% Selected configuration: experiments/external_dataset/results/D1-dev-tuning/summary.json.
The system returns the top three results, each an
article with an excerpt and its source document. For BM25 or dense retrieval
we quote the chunk with the highest score for that method. For hybrid retrieval
with $0<\lambda<1$, we apply each method's article-level min--max transform to
every chunk in the article and quote the chunk with the largest weighted sum.
Values below the article-level minimum may be negative. A method with zero
range contributes zero. At either endpoint we use the corresponding single
method. Ties are broken by ascending article or chunk identifier.
For HCMUS we group articles by topic within the selected cohort edition.
A general search shows the newer edition first. Two editions can cover the same topic
and still prescribe different rules. For general graduation eligibility queries,
the HCMUS interface checks each edition for a unique numbered clause that
introduces the eligibility list. If found, it promotes that article and
opens the cited list directly. Other questions use the usual ranking.

\textbf{Low-confidence warning.} Let $c^\star$ be the best chunk of the
top-ranked article before topic grouping in the original evaluation.
For ordinary ranked searches within a cohort or training mode, the demo uses the cosine
of the first excerpt within that selection. If $\mathrm{sim}(q,c^\star)<\theta$, the results are still
shown but are labelled as low confidence, and the student is advised to contact
the academic affairs office. Per-query min--max normalisation discards
absolute score levels, and the two methods may rank different articles first.
We therefore use the cosine for the warning.
% Selected threshold and coverage rule: experiments/external_dataset/results/D1-dev-tuning/summary.json.
We selected $\theta=0.6201629042625427$ on D1-dev to retain at least 80\% coverage.
We evaluated that operating point on D1-test and D2-test. D2 includes
questions outside the scope of the university regulations, which let us
measure how often the fixed threshold marks those results confident
(Section~\ref{sec:d2-results}). The system displays excerpts directly
from the indexed regulation texts.
